\begin{table}[htbp]
    \centering
    \caption{Cross-source equivalence within task type (TOST, relaxed margins, $\alpha=0.05$). Each cell reports equivalent source-pairs over all available source-pairs.}
    \setlength{\tabcolsep}{4pt}
    \begin{tabular}{lccc}
    \hline
    Metric & Margin & Single-hop Attribute Query & Entity Recognition \\
    \hline
    Question Length & $\pm5$ & 1/1 (1.00) & 3/3 (1.00) \\
    Answer Length & $\pm10$ & 0/1 (0.00) & 1/3 (0.33) \\
    Text-alone Yes Ratio & $\pm0.2$ & 1/1 (1.00) & 3/3 (1.00) \\
    External-evidence Yes Ratio & $\pm0.2$ & 1/1 (1.00) & 1/3 (0.33) \\
    RTE & $\pm0.1$ & 0/1 (0.00) & 0/3 (0.00) \\
    LJ & $\pm0.1$ & 0/1 (0.00) & 0/3 (0.00) \\
    OK State Ratio & $\pm0.2$ & 1/1 (1.00) & 0/3 (0.00) \\
    \hline
    \end{tabular}
    \label{tab:cross_source_consistency_one_table_relaxed}
\end{table}

\paragraph{Interpretation.}
Entries report TOST equivalence pass counts under relaxed practical margins; higher pass rates indicate stronger cross-source consistency for the metric.

\paragraph{Conclusion.}
With relaxed margins, structural properties (especially question length) show strong consistency, while process/performance metrics (RTE, LJ, and partly OK ratio) still exhibit notable source sensitivity, particularly in Entity Recognition.
